\documentclass[11pt]{article}

\usepackage[T1]{fontenc}
\usepackage{amsmath,amssymb,amsthm}
\usepackage{booktabs}
\usepackage[hidelinks]{hyperref}
\usepackage[margin=1in]{geometry}
\setlength{\emergencystretch}{3em}
\hypersetup{
  pdftitle={Length Authority and Representation Canonicality at an SQIsign Verification Boundary},
  pdfauthor={Dana Edwards},
  pdfsubject={Bounded implementation-security study of SQIsign verification interfaces},
  pdfkeywords={SQIsign, parser safety, signature representation, AddressSanitizer, Lean}
}

\newtheorem{proposition}{Proposition}
\newtheorem{observation}{Observation}
\newcommand{\code}[1]{\texttt{#1}}

\title{Length Authority and Representation Canonicality\\
at an SQIsign Verification Boundary}
\author{Dana Edwards}
\date{Version 1.0 preprint, 12 August 2026}

\begin{document}

\maketitle

\begin{center}
\fbox{\parbox{0.91\linewidth}{\centering
\textbf{PREPRINT --- NOT PEER REVIEWED.}
The findings concern one pinned public implementation revision and a bounded
x86 test matrix.  SQIsign maintainers confirmed the intended API contracts,
reported independently developed fixes in their development tree, reproduced
the affected behavior, and authorized immediate publication.}}
\end{center}

\begin{abstract}
We study length authority and representation canonicality at the public
verification boundary of the SQIsign implementation at revision
\code{dd133d7aca576c361a270c8e6434832535b42ecc}.  On 11 August 2026 the
official remote \code{main} branch resolved to that revision, and a fresh
same-machine replay exercised a six-authority matrix comprising parameter
levels 1, 3, and 5 under the reference and Broadwell x86 implementations.
Across both public verification
paths, all 2,656 tested undersized inputs produced AddressSanitizer findings
before a tested boundary repair; the repair converted all tested cases to
clean pre-decode rejections.  The detached verifier also accepted all 1,200
tested representations formed by appending one or sixteen declared bytes to
600 valid fixed-size KAT signatures.  An exact-length guard rejected all such
tested suffix aliases while preserving all 600 exact KAT signatures.  Generic
Lean theorems identify the structural boundary:
accepted decoder fibers are injective exactly when accepted encodings are
normalization fixed points, while the ordered x-only observable
$(x(R),x(S),x(R-S))$ determines $(R,S)$ only up to simultaneous negation.
The results are bounded implementation and representation findings, not a
forgery, key recovery, remote-code-execution result, deployment exploit, or
defect in the mathematical security proof.  The fresh replay packet commits
the runner, environment, complete per-case records, representative raw logs,
and a fail-closed content-addressed receipt.  Following private disclosure,
the maintainers reported that they had independently implemented equivalent
boundary checks approximately one month earlier.  They reproduced the affected
behavior by reverting those checks and reported clean rejection without
AddressSanitizer findings after restoring them.  They also confirmed that exact
equality is the intended detached-signature contract and authorized immediate
publication.  Their confirmation is independent maintainer evidence, but no
maintainer raw logs or public fixing commit were available at this preprint's
cutoff.
\end{abstract}

\paragraph{Keywords.}
SQIsign; parser safety; signature representation; canonical encoding;
AddressSanitizer; formal verification; Lean.

\section{Introduction}

A verification function receives bytes together with an asserted extent.  A
fixed-size decoder may safely consume those bytes only after the extent has
been checked.  This apparently mundane ordering constraint is a security
boundary: once a bare pointer enters a decoder, the evidence that the required
bytes exist is no longer present in the decoder's interface.

SQIsign is a compact post-quantum signature scheme built from supersingular
isogenies and quaternion arithmetic~\cite{sqisign-paper}.  This paper develops
a bounded case study for its public implementation
repository~\cite{sqisign-implementation}.  It separates three questions that
are easy to conflate:

\begin{enumerate}
  \item Does every public verification path reject an undersized input before
        fixed-size decoding?
  \item Does a detached signature's declared length participate in the
        accepted representation, or is only a fixed prefix consumed?
  \item When two byte strings decode to the same mathematical data, can both
        reach an accepting verifier state?
\end{enumerate}

The first property is false throughout the complete tested undersized x86
matrix at the pinned revision.  The second question yields accepted suffix
aliases at the detached pointer/length interface.  The third requires a more
careful answer: a generic formal theorem shows that byte uniqueness depends on
an accepted-normal-form property beyond mathematical object uniqueness.

Our contributions are deliberately narrow.

\begin{itemize}
  \item We record a complete tested-domain falsification of pre-decode length
        authority for two public APIs across six x86 authorities.
  \item We record accepted overlong detached representations and distinguish
        this API-level fact from a strong-unforgeability conclusion.
  \item We evaluate a minimal boundary repair against undersized inputs,
        exact KAT signatures, and two suffix lengths, and record subsequent
        maintainer confirmation of both the intended contracts and an
        independently developed development-tree fix.
  \item We connect the implementation findings to generic normalization and
        x-only quotient formal sources while retaining their distinct build
        and provenance status, without claiming a formal refinement theorem
        for production SQIsign.
\end{itemize}

\section{Authority, evidence, and claim discipline}

\subsection{Frozen implementation authority}

The audited authority is the repository \code{SQISign/the-sqisign} at the
40-hex revision
\begin{center}
\code{dd133d7aca576c361a270c8e6434832535b42ecc}.
\end{center}
The recorded matrix covers parameter levels 1, 3, and 5, each under the
reference and Broadwell x86 implementation.  We make no claim about another
revision, architecture, backend, API wrapper, or deployment.  Live remote
queries on 11 and 12 August 2026 resolved official \code{main} to this exact
commit.  These time-stamped equalities are kept separate from the frozen
authority and do not conflict with the maintainers' statement that the fixes
were already present in a separate development tree.

\subsection{Evidence classes}

All quantitative claims in this preprint use a fresh end-to-end run from a
clean official checkout rather than the exploratory discovery records.
The fresh packet contains the exact runner and harnesses, GCC/ASan environment,
all per-case records for the complete matrix, representative sanitizer and
safe-return logs, and hashes for every generated artifact.  Its normal and
optimized validators agree.  AddressSanitizer is the dynamic memory-error
detector described by Serebryany et al.~\cite{asan-paper}.  The packet closes
the local-replay and dated-source gaps for the pinned revision; it is not an
independent third-party reproduction.

On 12 August 2026, following private disclosure, the SQIsign team supplied a
written maintainer confirmation.  They stated that the two checks had been
developed independently approximately one month before the report, reproduced
the affected behavior after temporarily undoing those checks, and obtained
clean rejections without AddressSanitizer findings after restoring them.  This
is an independent maintainer reproduction and contract confirmation, although
the project does not possess their raw logs or development-tree commit.  A
sanitized correspondence summary and disclosure timeline accompany this
preprint; the private email address and full message headers are not published.

No quantitative result in this preprint rests solely on a conversational or
user-reported count.  The evidence mapping and nearest nonclaim for the six
headline claims appear in \path{claims_evidence.json}.

\section{Length authority before fixed-size decoding}

Let $s_1=148$, $s_3=224$, and $s_5=292$ denote the signature byte sizes in the
three recorded parameter profiles.  For each of two implementations and two
public verification paths, the campaign called the API at every declared
length $0,\ldots,s_i-1$.  The zero-length case used a one-byte non-null
allocation, separating the length contract from null-pointer behavior.  The
number of calls was therefore
\[
  2\text{ APIs}\cdot 2\text{ implementations}\cdot
  (148+224+292)=2{,}656.
\]

The two paths were the NIST-style \path{crypto_sign_open} interface and the
detached \path{sqisign_verify} interface.  At the pinned revision both paths
allowed a bare signature pointer to reach fixed-size decoding before the
provided length authorized the read.

\begin{table}[ht]
\centering
\caption{Complete recorded undersized-input matrix.}\label{tab:undersized}
\begin{tabular}{lrrr}
\toprule
Variant & Calls & ASan findings & Clean pre-decode rejections \\
\midrule
Pinned implementation & 2,656 & 2,656 & 0 \\
Tested repair          & 2,656 & 0     & 2,656 \\
\bottomrule
\end{tabular}
\end{table}

\begin{observation}[Bounded length-authority defect]
For every input in the recorded undersized x86 matrix, the pinned
implementation produced an AddressSanitizer finding rather than a clean
pre-decode rejection.  This falsifies the tested property that each public API
rejects every undersized input before fixed-size signature decoding.
\end{observation}

This observation is an implementation-safety result under the recorded build
and sanitizer configuration.  It does not by itself establish remote code
execution, confidentiality loss, or exploitability in any named deployment.

\section{Detached suffix aliases}

The detached verifier receives a signature pointer and a \code{siglen}
argument.  In the pinned implementation, the recorded source audit found that
\code{siglen} did not govern entry into the fixed-size signature decoder.
Consequently, the campaign evaluated each of 600 valid detached KAT
signatures in three forms: exact length, one appended declared byte, and
sixteen appended declared bytes.

The maintainers subsequently confirmed that exact equality is the intended
contract for detached \code{siglen}.  By contrast, the signed-message open path
correctly admits \code{smlen >= SIGNATURE\_BYTES}, because bytes following the
fixed-size signature are the message.  Its missing boundary condition was only
the lower bound \code{smlen < SIGNATURE\_BYTES}.

\begin{table}[ht]
\centering
\caption{Recorded detached-signature results.}\label{tab:trailing}
\begin{tabular}{lrr}
\toprule
Representation & Pinned accepted & Repaired accepted \\
\midrule
Exact fixed-size signature       & 600/600   & 600/600 \\
Exact signature plus one byte    & 600/600   & 0/600 \\
Exact signature plus sixteen bytes & 600/600 & 0/600 \\
\bottomrule
\end{tabular}
\end{table}

At the C API boundary, these are distinct declared byte representations with a
common valid fixed-size prefix, message, and verification result.  We call
them \emph{suffix aliases}.  This terminology avoids silently assuming that a
formal security experiment defines the signature domain as arbitrary
length-delimited strings.  If the formal domain contains only fixed-size
signatures, the extra suffix may lie outside the mathematical signature rather
than constitute a second signature in that experiment.  Therefore
Table~\ref{tab:trailing} is not, by itself, a strong-unforgeability theorem or
an existential forgery.

\section{Boundary repair and maintainer confirmation}

Our tested patch inserts two dominating guards before signature decoding.
For the signed-message path it rejects when
\code{smlen < SIGNATURE\_BYTES}, setting the output message length to zero.
For detached verification it rejects unless
\code{siglen == SIGNATURE\_BYTES}.

The recorded repaired matrix supplies three complementary checks:

\begin{enumerate}
  \item all 2,656 undersized calls became clean rejections with no recorded
        AddressSanitizer finding;
  \item all 600 exact KAT signatures remained accepted; and
  \item all 1,200 tested suffix aliases were rejected.
\end{enumerate}

This is closure over the complete tested x86 matrix, not a proof about all C
callers or all parser defects.  The maintainers confirmed that these are the
intended policies.  They also reported that their team had independently added
equivalent pre-decode checks approximately one month before receiving this
report, and that improved regression tests would accompany the Round~3
release.  Because their development-tree source was not public at the preprint
cutoff, we do not claim our patch is byte-identical to theirs or identify a
public fixing commit.

\section{Why decoder aliasing is not enough}

Suppose two byte strings decode to the same internal object.  This decoder
alias alone does not establish an accepted representation collision: an
independent algebraic check may reject the common decoded object in both
cases.  Conversely, the suffix experiment in Section~4 directly establishes
accepted aliases at the detached pointer/length interface because the two
distinct declared byte strings receive the same accepting verdict.

The distinction is therefore between equality in a decoder fiber and
injectivity of the decoder restricted to accepted inputs.  The next section
formalizes the exact extra condition needed to pass from the former to the
latter.  We make no universal claim about fixed-length SQIsign serialization.

\section{Structural canonicality theorems}

The accompanying Lean development separates mathematical fibers from byte
normal forms.

\subsection{Normalization fixed points}

Let $D:B\to O$ decode byte representations, $E:O\to B$ encode objects, and
$V:O\to\{\mathsf{true},\mathsf{false}\}$ verify objects.  Define
\[
  N(b)=E(D(b)), \qquad A(b)=V(D(b)).
\]
Assume $D(E(o))=o$ for every object $o$.

\begin{proposition}[Accepted-fiber normalization criterion]
The decoder is injective on accepted fibers if and only if every accepted byte
string is a fixed point of $N$.  Moreover, any accepted $b$ with $N(b)\ne b$
immediately yields two distinct accepted representations, $b$ and $N(b)$,
with equal decoded objects.
\end{proposition}

\begin{proof}[Proof idea]
Normalization preserves decoded objects by the section law and therefore
preserves acceptance.  Accepted-fiber injectivity forces $N(b)=b$.
Conversely, two accepted bytes in one decoder fiber have one common normal
form; if both are fixed points, they are equal.  The frozen Lean~4
source~\cite{lean4-paper}
\code{CanonicalNormalization.lean} contains this generic proof without
placeholders.  Its exact frozen source hash was freshly rebuilt with Lean
4.19.0 and the recorded Mathlib revision on 11 August 2026.  The fresh axiom
audit reports no axioms for the two normalization theorems used here.
\end{proof}

This proposition captures the inner ground for canonical acceptance: accepted
bytes must already be the selected normal form.  It does not instantiate the
pinned SQIsign decoder or verifier.

\subsection{The ordered x-only quotient}

For points on an affine Weierstrass curve over a field, define
\[
  X(R,S)=\bigl(x(R),x(S),x(R-S)\bigr),
\]
using a separate value for the point at infinity.  The classical elliptic
$x$-line is a sign quotient; related quotient geometry appears
in~\cite{costello-smith-2017,kohel-2016}.  The accompanying Lean development
checks the exact ordered-pair statement:
\[
  X(R',S')=X(R,S)
  \quad\Longleftrightarrow\quad
  (R',S')=(R,S)\ \text{or}\ (R',S')=(-R,-S).
\]
Thus the induced map from the simultaneous-sign quotient is injective.  The
third coordinate synchronizes the otherwise independent endpoint signs via
the identities $S=R-(R-S)$ and $R=(R-S)+S$.

This is a mathematical object theorem.  Byte uniqueness additionally needs a
production refinement showing that every accepted serialization is the chosen
orbit normal form.  No such refinement theorem is claimed here.

\section{Security interpretation and limitations}

The evidence supports two distinct bounded conclusions.  First, supplied
lengths did not dominate fixed-size decoding for either tested verification
path.  Second, detached verification accepted tested suffix aliases at its
pointer/length boundary.  Neither conclusion depends on treating the generic
x-only quotient theorem as a production verifier theorem.

The nearest likely overclaims are expressly excluded:

\begin{itemize}
  \item no invalid signature for a new message was produced;
  \item no secret key was recovered;
  \item no remote-code-execution or deployment-specific exploitability claim
        is made;
  \item no flaw in the mathematical SQIsign security proof is established;
  \item no universal serialization-canonicality theorem is established;
  \item no EUF-CMA or SUF-CMA conclusion follows without fixing and analyzing
        the exact formal representation domain;
  \item neither our tested patch nor the reported development-tree fix is
        claimed to eliminate every parser defect; and
  \item the maintainers' development-tree fix is reported correspondence
        evidence, not a source-level audit of a public fixing commit.
\end{itemize}

The complete matrix was rerun by the published end-to-end runner, with complete
case records and representative raw logs bound by the fresh receipt.  The
maintainer response adds independent human reproduction and contract
confirmation, but not their raw logs or development-tree source.

\section{Coordinated-disclosure timeline}

The report was sent privately to the SQIsign team on 11 August 2026.  On
12 August, the team acknowledged it, confirmed both intended length contracts,
reported the independently developed development-tree fixes and reproduction
described above, and stated that SQIsign is a reference implementation without
a responsible-disclosure policy.  The team explicitly authorized immediate
publication, including publication as a public issue, and said the fixes would
ship in the forthcoming Round~3 release.

At the preprint cutoff, the public \code{main} branch still resolved to the
affected revision.  We therefore distinguish three states: the frozen public
revision tested here; the maintainer-reported development-tree repair; and the
future public fixing revision.  This chronology avoids both implying that the
report caused a fix developed earlier and implying that an unpublished repair
was already present in the audited public source.

\section{Reproducibility map}

The release directory is designed to stand alone as a GitHub or Zenodo
artifact and is available at
\url{https://github.com/TheDarkLightX/sqisign-verification-boundary}.
Its public artifact map is as follows:

\begin{sloppypar}
\begin{itemize}
  \item tested source patch:
    \path{patches/sqisign-dd133d7-h7-length-guards.patch};
  \item fresh full-matrix runner, case records, representative raw logs, and
    receipt: \path{research/h7_completion_2026_08_11/};
  \item formal normalization theorem:
    \path{research/h7_completion_2026_08_11/lean/IsogenyCrypto/H7/CanonicalNormalization.lean};
  \item formal x-only fiber and quotient theorems:
    \path{research/h7_completion_2026_08_11/lean/IsogenyCrypto/H7/XCoordinateFiber.lean}
    and
    \path{research/h7_completion_2026_08_11/lean/IsogenyCrypto/H7/XOnlyQuotient.lean};
  \item exact-source Lean 4.19.0 build and axiom receipt:
    \path{research/h7_completion_2026_08_11/lean/validation_receipt.json};
  \item sanitized maintainer confirmation and disclosure chronology:
    \path{MAINTAINER_CONFIRMATION.md} and
    \path{DISCLOSURE_TIMELINE.md}.
\end{itemize}
\end{sloppypar}

Exact git refs, evidence classes, selected hashes, and nonclaims for the five
headline claims are machine-readable in \path{claims_evidence.json}.

\section*{Acknowledgments and AI-use disclosure}

The author thanks D\'ecio and the SQIsign team for their
prompt response, independent reproduction, contract clarification, and
disclosure guidance.  OpenAI Codex agents assisted with exploratory analysis,
test and validation code, evidence organization, and drafting.  Their outputs
were treated as untrusted candidate material and checked against frozen source,
replayable experiments, and machine-checked proof artifacts.  Dana Edwards
reviewed and approved the claims, experiments, limitations, and final text and
takes responsibility for the work.  No AI system is an author.

{\small
\bibliographystyle{plain}
\bibliography{references}
}

\end{document}
